\subsection{The 3-Market Opportunity Ledger}

\begin{table*}[t]
\centering
\small
\begin{tabularx}{\textwidth}{>{\RaggedRight\bfseries}p{3.2cm} >{\RaggedRight}X >{\RaggedRight}X >{\RaggedRight}X}
\toprule
Decision Variable & D2C Engineer & College Partner & Enterprise Team \\
\midrule
Who Experiences Pain? & Software developer & Faculty and students & Engineering team \\
\addlinespace
Who Controls Payment? & Individual developer (savings/credit) & Principal / Treasurer (tuition/grant) & VP Engineering (R\&D budget) \\
\addlinespace
Dominant Job & Defend code; master production systems & Elevate curriculum; accreditation & Eliminate release defects; hit deadlines \\
\addlinespace
Primary Substitute & Docs; senior peer; self-study & Internal faculty; online MOOCs & Internal mentorship; external agency \\
\addlinespace
Primary Anxiety & Wasting personal money/time & Academic calendar clash; faculty revolt & Infosec breach; code leakage \\
\addlinespace
First Testable Proof & Paid 2-day architectural studio & Paid 2-week winter laboratory & Paid evaluation diagnostic pilot \\
\addlinespace
Founder Bottleneck & Time spent reviewing code line-by-line & Academic committee negotiations & Long sales cycles and infosec reviews \\
\bottomrule
\end{tabularx}
\caption{Cross-market comparison ledger across Bayes Institute's 3 candidate ecosystems.}
\label{tab:ch6_three_market_ledger}
\end{table*}

\subsection{Which Segment to Interview First?}

A solo founder with 15 to 20 weekly hours cannot execute discovery across all 3 markets simultaneously. Spreading scarce attention across individual developers, university deans, and corporate procurement directors guarantees shallow insights and zero traction.

\noindent\textbf{Recommended Discovery Sequence:}
\begin{enumerate}[leftmargin=1.2em,itemsep=0.2ex]
  \item \textbf{Cycle 1 (Weeks 1 to 4): D2C Software Engineers.} Fastest discovery velocity; zero procurement friction; immediate feedback on technical curriculum depth.
  \item \textbf{Cycle 2 (Weeks 5 to 8): Technical Engineering Managers (B2B).} High contract values; validates whether enterprise engineering leaders will fund the same capability for teams.
  \item \textbf{Cycle 3 (Subsequent Term): Engineering Colleges.} Defer until a proven, repeatable instructional and lab assessment architecture exists.
\end{enumerate}

\subsection{What Real Evidence Changes the Curriculum Design}

True customer understanding alters what you build:
\begin{itemize}[leftmargin=1.2em,itemsep=0.2ex]
  \item If field discovery reveals engineers understand theory but struggle with latency and compute costs: strip out introductory math and build an intensive benchmarking studio.
  \item If discovery shows engineers cannot commit 8 hours on Saturdays due to family duties: redesign the curriculum as an asynchronous, code-review-driven sprint with two 90-minute evening check-ins.
  \item If discovery shows that prospective learners want degree credentials for overseas immigration: explicitly disqualify that segment and refuse to build a diluted mock-university program.
\end{itemize}

\subsection{Testing Capability Demand Without Career Promises}

Bayes Institute refuses on principle to deploy deceptive placement marketing or promise guaranteed salary hikes. How can it test demand for pure capability?

Administer an \textbf{Independent Systems Audit Challenge}:
\begin{itemize}[leftmargin=1.2em,itemsep=0.2ex]
  \item Provide a public, open-source repository containing an AI evaluation pipeline with 4 deliberately injected, realistic production defects (data leakage, metric gaming, latency bottleneck, prompt injection vulnerability).
  \item Invite qualified engineers to submit pull requests identifying and repairing the flaws within 48 hours.
  \item Provide detailed, rigorous, line-by-line code reviews on the top 15 submissions.
\end{itemize}
Observe how participants respond: Do they value the depth of critique? Do they request continued, structured feedback on their live work? Are they willing to pay a fair fee for an intensive 6-week studio? This tests demand for \textbf{genuine capability transformation} rather than certificate hunting.

\subsection{The Bayes Institutional Evidence Dashboard}

Track 4 distinct empirical metrics:
\begin{enumerate}[leftmargin=1.2em,itemsep=0.2ex]
  \item \textbf{Problem Urgency Score:} Documented hours lost and deadlines delayed across interviewees.
  \item \textbf{Behavioral Commitment Rate:} Percentage of qualified applicants who submit demanding technical readiness diagnostics.
  \item \textbf{Transformation Ratio:} Baseline diagnostic score vs. blind endline assessment score on unseen production architectures.
  \item \textbf{Founder Operational Load:} Total instructor and code-review hours required per completed student outcome.
\end{enumerate}

\subsection{The Uncomfortable Decision Bayes Must Be Ready to Make}

Suppose 60 days of disciplined field research reveal the following empirical facts:
\begin{itemize}[leftmargin=1.2em,itemsep=0.2ex]
  \item Individual engineers praise the curriculum, but 85\% refuse to pay more than ₹5,000 from personal savings because their employers reimburse technical books.
  \item However, 4 engineering managers immediately offer to pay ₹2,50,000 from team discretionary budgets to have Bayes audit their evaluation harnesses and upskill their 6-person teams.
\end{itemize}

The founder must possess the strategic discipline to \textbf{pivot the commercial entry point from D2C consumer courses to B2B enterprise engineering studios}. The enduring institutional mission: developing human capability: remains unchanged; the commercial entry vehicle adapts to market reality.

\subsection{Customer Knowledge as an Enduring Asset}

In a seed-stage startup, customer understanding lives inside the founder's head. In an enduring institution, customer understanding must be codified into an \textbf{institutional asset that survives the founder}:
\begin{itemize}[leftmargin=1.2em,itemsep=0.2ex]
  \item A living \textbf{Cognitive Defect Taxonomy}: documented common misconceptions and technical debugging hurdles encountered by learners across cohorts.
  \item A standardized \textbf{Rubric of Systems Independence}: objective criteria used to evaluate whether an engineer can safely ship code to production.
  \item Secure, privacy-preserving case archives of real-world production architectures.
\end{itemize}
Codifying customer insight transforms individual founder intuition into institutional pedagogical capability.
